% ===================================================================
%  16. TAULA — Denda handia. — Bazarra. — Jostailuak.
%  Traduction 2026 d'apres texto/io, controlee sur texto/fr, texto/en
%  et texto/es. Consigne : voir texto/eu/10-taula-01.tex.
%
%  Le tableau d'astronomie (bloc 15-1) porte des renvois a LETTRES,
%  (a) a (f), graves sur le tableau lui-meme : ils se rangent sous le
%  numero 85, comme le dit gravuri/literi.json. Ce bloc est aussi le
%  troisieme ecart declare de renvoji.py.
%
%  DERNIER TABLEAU DE LA COLONNE BASQUE. Bilan des seize : voir le
%  message de commit, qui porte le compte des inversions et ce qu'il
%  faut en conclure.
% ===================================================================

%%K t16-apar-1 apar
\VUsaut{4.0mm}
\VUtitre{15.0pt}{60}{16. TAULA}
\VUsaut{2.0mm}
\VUpk{13.2pt}{40}{Denda handia. --- Bazarra.}
\VUpk{13.2pt}{40}{Jostailuak.}
\VUsaut{3.4mm}

%%K t16-01-1 p
1. --- Urte berria hurbiltzen ari da. Denda handiek prestatu dituzte beren
\VUgras{jostailuen} \textsuperscript{(1)} eta urte berriko opariren erakusketak.

%%K t16-01-2 p
Aitonari eskatu nion bazarrera eramateko erakusketa eder hori ikustera, eta
onartu zuen.

%%K t16-02-1 p
2. --- Lehenik arma-panel ederren aurrean gelditu ginen, non
\VUgras{eskopeta}, \VUgras{txanoa}, \VUgras{sablea}, \VUgras{ezpata-uhala} eta
\VUgras{sorbaldako} pare bat baitzeuden, edo bestela \VUgras{kaskoa},
\VUgras{kirasa} \textsuperscript{(3)} eta \VUgras{pistola}
\textsuperscript{(4)}. Horrek guztiak askoz gehiago interesatzen ninduen
\VUgras{argi-linterna magikoek} \textsuperscript{(2)} baino.

%%K t16-tit-1 sub
\VUsaut{3.0mm}
\VUpk{10.2pt}{40}{Ikuskizun ederrak.}
\VUsaut{2.6mm}

%%K t16-03-1 p
3. --- Sail berean \VUgras{zaindari} bat \textsuperscript{(5)} zegoen bere
\VUgras{etxolan}; kartoizko animalia-erakusketa oso bat zaintzen zuela zirudien.
Zenbat animalia zeuden han! \VUgras{Loroa} \textsuperscript{(6)} bere makilan,
\VUgras{errinozeroa} \textsuperscript{(7)} adarra sudurrean, \VUgras{ipar oreina}
\textsuperscript{(8)}, \VUgras{ostruka} \textsuperscript{(9)},
\VUgras{tximinoa} \textsuperscript{(10)} intxaur bat karraskatzen,
\VUgras{azeria} \textsuperscript{(11)} oilo bat eramaten, \VUgras{krokodiloa}
\textsuperscript{(12)} aho erraldoia zabalik, \VUgras{gamelua}
\textsuperscript{(13)}, \VUgras{galgo} lirain bat \textsuperscript{(14)},
\VUgras{pantera} \textsuperscript{(15)}, eta gero ezagutzen ez nituen bi animalia,
zeinak, esan zidatenez, \VUgras{erbinudea} \textsuperscript{(16)} eta
\VUgras{hiena} \textsuperscript{(17)} baitziren. Bazegoen oraindik
\VUgras{lehoinabarra} \textsuperscript{(18)} eta \VUgras{otsoa}
\textsuperscript{(21)}; ondo-ondoan \VUgras{arratoi} politak
\textsuperscript{(19)} eta gomazko \VUgras{sagu} txikitxoak
\textsuperscript{(20)}. Azkenik, \VUgras{hipopotamoak} \textsuperscript{(22)}
eta \VUgras{dromedario} konkordunak \textsuperscript{(23)} osatzen zuten bilduma.
Han egongo nintzatekeen orduak eta orduak, ondoan \VUgras{beruzko soldadutxoz}
\textsuperscript{(29)} betetako kanpaleku bikain bat egon ez balitz, zeinak nire
arreta erakarri baitzuen.

%%K t16-tit-2 sub
\VUsaut{4.0mm}
\VUpk{10.2pt}{40}{Soldaduak eta gerra.}
\VUsaut{2.8mm}

%%K t16-04-1 p
4. --- Aitonak, zeina \VUgras{elbarria} baita \textsuperscript{(24)} eta zerbitzua
utzi behar izan baitzuen \VUgras{zauri} bat zela eta --- zeinak
\VUgras{gerra-domina} ekarri baitzion ---, gogoko du parte hartu zituen
\VUgras{kanpainak} kontatzea. Gogotik azaldu zizkidan miniaturazko armada txiki
haren mugimendu desberdinak. Bazarra ikusten ari zen benetako soldadu talde bat ---
\VUgras{zapalaria} \textsuperscript{(25)}, \VUgras{mendiko tiratzailea}
\textsuperscript{(26)} bere \VUgras{txapelarekin}, infanteriako
\VUgras{sarjentu nagusia} \textsuperscript{(27)} eta artilleriako
\VUgras{sarjentua} \textsuperscript{(28)} urrezko \VUgras{galoiarekin} besoan ---
gure ondoan gelditu zen azalpenak entzutera.

%%K t16-05-1 p
5. --- Aitonak, entzule hain agurgarriaz oso harro, bat-batean bataila-plan oso bat
inprobisatu zuen.

%%K t16-05-2 p
Gotorlekua, bere \VUgras{harresi} eta \VUgras{lubanarroekin}, \VUgras{etsai-armadak}
setiatu zuen, zeinak \VUgras{bonbardaketa} luze baten ondoren
\VUgras{oldarrez} hartzeko prestatzen baitzen. Baina \VUgras{setiatuek}, goseak
hertsatuta, \VUgras{irteera} bat egin zuten \VUgras{setiatzaileen}
\VUgras{kanpalekuaren} kontra. \VUgras{Erasoa} \VUgras{borroka} gogorrean
\VUgras{denden} inguruan uxatuta, setiatzaile \VUgras{garaileek} beren
\VUgras{eskuadroiak} bidali zituzten \VUgras{jazarpenera} eraso egin zuten
\VUgras{garaituen} atzetik. Haiek \VUgras{atzera egin zuten}, \VUgras{gudu-zelaia}
\VUgras{hildakoz} eta \VUgras{zaurituz} josita utzita. \VUgras{Osasun-langileek}
zaurituak \VUgras{kanpainako ospitalera} eraman zituzten,
\VUgras{kirurgialariak} senda zitzan.

%%K t16-05-3 p
\VUgras{Batailoi} batzuen erresistentzia heroikoa gorabehera --- zeinak
\VUgras{koadroan} lerrokatu baitziren ---, \VUgras{porrota} erabatekoa izan zen;
armadaren hondarrak \VUgras{ihes egin zuen}, \VUgras{preso} asko utzita. Etsaia
bere \VUgras{garaipena} gotorlekua hartuz koroatzeko prestatzen zen. Agian
kanpainaren amaiera izango zen, eta \VUgras{su-etenaren} ondoren
\VUgras{bake-ituna} sinatu ahal izango zen.

%%K t16-tit-3 sub
\VUsaut{4.4mm}
\VUpk{10.2pt}{40}{Urte berriko opariak.}
\VUsaut{2.8mm}

%%K t16-06-1 p
6. --- Soldadu gazteek, beteranoaren kontakizun bizia entzunez barre egiten
zutenek, galdera batzuk gehiago egin zizkioten. Niri dagokidanez, mahaia inguratu
nuen \VUgras{balezta} eder bati \textsuperscript{(32)} eta \VUgras{arku} bati
\textsuperscript{(33)} begiratzeko, bere \VUgras{geziarekin}
\textsuperscript{(31)}. Han beste animalia-bilduma bat aurkitu nuen: bi
\VUgras{txori kantari} \textsuperscript{(34)}, \VUgras{urretxindor} eta
\VUgras{txinbo} mekanikoak, ahal zuten guztia kantatzen, eta izaki bitxi
batzuen lerroa --- \VUgras{saguzarra} \textsuperscript{(35)},
\VUgras{boa} \textsuperscript{(36)}, \VUgras{dortoka}
\textsuperscript{(37)}, \VUgras{itsas txakurra} \textsuperscript{(38)},
\VUgras{balea} \textsuperscript{(39)} eta \VUgras{marrazoa}
\textsuperscript{(40)}, mintzez eginak.

%%K t16-07-1 p
7. --- Ez nien luze begiratu, amesten nuena ikusi bainuen: \VUgras{txotxongilo-antzoki}
bikain bat \textsuperscript{(41)}, \VUgras{dekoratu} ederrekin eta
\VUgras{oihal} gorri xarmangarri batekin. \VUgras{Eszenatokian} bi aktorek
\VUgras{papera} antzezten zutela zirudien \VUgras{antzezlan} interesgarrienean, izan
ere kartoizko \VUgras{ikusle} txikiak palkoetan --- eta baita \VUgras{platean} eta
\VUgras{galerian} \VUgras{zuzendariaren} atzean eserita zeudenak ere --- gogotik
txalotzen ari baitziren.

%%K t16-07-2 p
Zoritxarrez, antzoki hura oso garestia zen.

%%K t16-08-1 p
8. --- Gainera zaila zen jostailuen artean aukeratzea, erakusleihoak mota guztietako
jostailuz beteta baitzeuden: \VUgras{Arlekina} \textsuperscript{(42)},
\VUgras{Polizinela} \textsuperscript{(43)}, \VUgras{Pierrot}
\textsuperscript{(44)}, hegoak astintzen zituzten \VUgras{hontzak}
\textsuperscript{(45)}, txistuka ari ziren \VUgras{birigarroak}
\textsuperscript{(46)}, zaunkaka ari ziren \VUgras{kanitxeak},
\VUgras{gramofonoa} \textsuperscript{(47)} eta \VUgras{buruhausteak}. Jostailu
horietako batek asko dibertitu ninduen: \VUgras{itsas gudu} bat
\textsuperscript{(48)} irudikatzen zuen, non \VUgras{ontzi} bati
\VUgras{piratek} eraso egiten baitzioten \VUgras{koko-palmondoz} landatutako
itsasertz baten ondoan.

%%K t16-noto-1 noto
\VUnotes{143.2mm}{%
\textit{(*) Ikus 6. taula.}%
}

%%K t16-09-1 p
9. --- Mirari horiek guztiak lasai-lasai begira nitzakeen, dendan jende gutxi
baitzegoen --- begiluze eskukada bat ozta-ozta, \VUgras{dragoi} bat
\textsuperscript{(49)} eta \VUgras{kobratzaile} bat \textsuperscript{(50)},
\VUgras{letra} bat aurkezten \VUgras{kutxa} nagusian \textsuperscript{(51)}.

%%K t16-10-1 p
10. --- Aitonari galdetu nion zertarako ziren \VUgras{kutxazainaren} atzeko
apaletako liburuak; esan zidan \VUgras{kontabilitate-liburuak} zirela.

%%K t16-tit-4 sub
\VUsaut{3.6mm}
\VUpk{10.2pt}{40}{Denda begiratzen.}
\VUsaut{2.8mm}

%%K t16-11-1 p
11. --- Jostailuen sailetik atera eta zaldi-tresnenera joan ginen
\VUgras{zelak} \textsuperscript{(53)}, \VUgras{estriboak}
\textsuperscript{(54)}, \VUgras{bridak} \textsuperscript{(55)},
\VUgras{ahoburdinak} \textsuperscript{(56)} eta \VUgras{zartailuak} ikustera.
\VUgras{Sail-buruak} \textsuperscript{(57)} aitonari argibideak ematen zizkion
bitartean, \VUgras{portzelanaren} eta \VUgras{beiraren} \textsuperscript{(58)}
erakusketa ikustera joan nintzen, non \VUgras{entsalada-ontzi} ederrak eta
\VUgras{te-ontzi} dotoreak \textsuperscript{(59)} baitzeuden.

%%K t16-12-1 p
12. --- Lehen solairura igotzeko, \VUgras{kortseen} \textsuperscript{(60)}
erakusleihoaren atzetik igaro ginen, galtzerdien sailaren ondotik, non
\VUgras{galtzontzilo} ederrak \textsuperscript{(63)} baitzeuden zabalduta. Ondoan
\VUgras{saltzaile} batek \textsuperscript{(61)} \VUgras{erosle} bati
\textsuperscript{(62)} laguntzen zion zetazko \VUgras{oihal} ederrak
\textsuperscript{(64)} aukeratzen.

%%K t16-12-2 p
Eskaileretatik igotzean, ohartu nintzen denda japoniar \VUgras{argiontziz}
\textsuperscript{(65)}, \VUgras{eguzkitakoz} \textsuperscript{(66)} eta
\VUgras{palmondo} erraldoiez \textsuperscript{(70)} apainduta zegoela.

%%K t16-13-1 p
13. --- Lehen solairuan ia ez zegoen ezer ikusteko: altzariak (*) baino ez,
\VUgras{kutxa gotorrak} \textsuperscript{(67)}, \VUgras{komodak}
\textsuperscript{(68)} eta \VUgras{haur-kotxeak} \textsuperscript{(69)}. Aitzitik,
dendaren ikuspegi orokorra oso \VUgras{interesgarria} zen.

%%K t16-14-1 p
14. --- Gure azpian musika-tresnak ikusten nituen: \VUgras{harpak}
\textsuperscript{(71)}, \VUgras{gitarrak} \textsuperscript{(72)},
\VUgras{mandolinak} \textsuperscript{(73)}, \VUgras{kornetak}
\textsuperscript{(74)}, \VUgras{klarineteak} \textsuperscript{(75)}, biolinak
beren \VUgras{arkuekin} \textsuperscript{(76)} eta \VUgras{danbor handia}
\textsuperscript{(77)} ere. Pixka bat urrunago, paper-mahaian, \VUgras{ikasle}
batek \textsuperscript{(78)} \VUgras{saltzailearekin} \textsuperscript{(79)}
paper-zorro baten prezioa eztabaidatzen zuen. Haien ondoan \VUgras{abezedario}
eder bat \textsuperscript{(80)} bereizten nuen.

%%K t16-14-2 p
Dendaren erdian \VUgras{Gabon-zuhaitz} garai bat \textsuperscript{(81)} zegoen,
kandelaz, apaingarriz eta mota guztietako jostailuz betea: \VUgras{zurezko
zalditxoez} \textsuperscript{(82)}, \VUgras{panpinez} \textsuperscript{(83)} eta
abarrez.

%%K t16-14-3 p
\VUgras{Lurrin-sailak} \textsuperscript{(84)}, bere \VUgras{hortz-uren} eta
\VUgras{lurrin} desberdinekin, usain oso atsegina zabaltzen zuen, guregana
igotzen zena.

%%K t16-tit-5 sub
\VUsaut{3.4mm}
\VUpk{10.2pt}{40}{Zerbait erosten dugu.}
\VUsaut{2.8mm}

%%K t16-15-1 p
15. --- Aitona aspertzen ari zenez, berriro jaitsi ginen eskailera handitik. Une
batez gelditu zen \VUgras{astronomia-taularen} \textsuperscript{(85)} aurrean,
zeinean, esan zidanez, hauek irudikatuta baitzeuden: \VUgras{kometak}
\textsuperscript{(a)}, \VUgras{ilargiaren eta eguzkiaren eklipseak}
\textsuperscript{(b)}, haize-arrosa \VUgras{munduaren lau alderdiekin}
\textsuperscript{(c)} \VUgras{(iparra, hegoa, ekialdea eta mendebaldea)}, bi
\VUgras{hemisferioen} mapa \textsuperscript{(d)}, \VUgras{planetak}
\textsuperscript{(e)} eta \VUgras{eguzki-sistema} \textsuperscript{(f)}. Nik ez
nuen ezer ulertu, eta askoz gehiago interesatzen ninduten ondotik igarotzen zen
\VUgras{mandatariaren} \textsuperscript{(86)} libre galoiduna eta nire ondoko
\VUgras{haizemaile} politak \textsuperscript{(87)}, \VUgras{diru-zorroen}
\textsuperscript{(88)} eta \VUgras{kartera-zorroen} \textsuperscript{(89)} ondoan.

%%K t16-16-1 p
16. --- Mahaian \VUgras{lumak} \textsuperscript{(90)} eta \VUgras{gerriko-belarriak}
\textsuperscript{(91)} ere miresten nituen, eta \VUgras{lore artifizial} politak
\textsuperscript{(94)}, saskietan dotore jarriak. \VUgras{Bioletak},
\VUgras{eguzkiloreak}, \VUgras{jazintak} \textsuperscript{(95)},
\VUgras{geranioak} \textsuperscript{(96)}, \VUgras{lirioak}
\textsuperscript{(97)} eta \VUgras{kapuxinak} \textsuperscript{(98)} oso ondo
imitatuta zeuden.

%%K t16-tit-6 sub
\VUsaut{4.4mm}
\VUpk{10.2pt}{40}{Aitonaren oparia.}
\VUsaut{2.8mm}

%%K t16-17-1 p
17. --- Aitonari eskatu nion \VUgras{hortzetako eskuiletako}
\textsuperscript{(92)} bat eros ziezadan, \VUgras{orratz-buruen}
\textsuperscript{(93)} ondoko kutxan zeudenetako bat. Onartu zuen, eta neuk joan
nintzen nire erosketa kutxan ordaintzera, zeinaren ondoan \VUgras{pakete} handi
bat \textsuperscript{(100)} lotzen baitzuten \VUgras{erosle} batentzat
\textsuperscript{(99)}.

%%K t16-18-1 p
18. --- Bat-batean, arte-\VUgras{brontzeen} \textsuperscript{(101)} ondoan zeuden
pertsonen artean asaldura txiki bat sortu zen. \VUgras{Lapur} bat
\textsuperscript{(102)} harrapatu berri zuen \VUgras{dendako detektibeak}
\textsuperscript{(103)} norbaiti lapurtzen saiatzen ari zen unean. Udaltzainaren
bila bidali zuten hura \VUgras{atxilotzeko}, eta ateratzen ari ginen unean bertan
errudunari komisariara eramaten zioten.
\VUsaut{6.0mm}
\VUfilet{22mm}
